% ============================================================
%  Chapter 1: Introduction
% ============================================================

\chapter{Introduction}
\label{ch:introduction}

\section{Motivation}

The design of artificial neural networks has historically been guided by
human intuition. From the earliest perceptron models of Rosenblatt
\citep{rosenblatt1958perceptron} to modern deep architectures
\citep{lecun2015deep}, researchers have manually specified the number of
layers, the connectivity patterns, the activation functions, and the
learning rules. This top-down approach has produced remarkable results,
but it has also imposed structural constraints that may not be optimal
for every task.

A central example of such a constraint is the layered architecture
itself. Since the backpropagation algorithm was popularized by
\citet{rumelhart1986learning}, nearly all successful neural networks
have been organized into discrete layers. This design is computationally
convenient—it allows efficient matrix multiplication and gradient
computation—but it also imposes a rigid separation between ``input,''
``hidden,'' and ``output'' units. In a biological brain, by contrast,
there are no layers. Neurons form a dense, recurrent graph in which any
neuron may connect to any other, and structure emerges through
developmental and evolutionary processes rather than explicit design.

This discrepancy raises a fundamental question: \emph{Can a neural
network discover its own structure from first principles, without a
designer specifying layers, connectivity, or even the number of nodes?}

\section{Background and Historical Context}

\subsection{The XOR Problem and the Limits of Linear Models}

The limitations of fixed, single-layer architectures were formalized by
\citet{minsky1969perceptrons}, who proved that a single-layer perceptron
cannot compute the XOR function, nor the parity function for any number
of inputs. This result, often referred to as the ``XOR problem,'' is not
merely a technical curiosity. It reflects a deeper mathematical fact:
linear threshold units can only separate linearly separable classes, and
parity is not linearly separable for any $n \geq 2$ \citep{minsky1969perceptrons}.

The XOR problem had a profound impact on the field. It contributed to
the first ``AI winter'' and demonstrated that architectural choices—in
particular, the presence or absence of hidden layers—are not merely
implementation details but determine the computational power of the
model. It was not until the rediscovery and popularization of
backpropagation by \citet{rumelhart1986learning} that multi-layer
networks became trainable, and the XOR problem was finally solved by
adding a single hidden layer.

\subsection{Evolutionary Approaches to Network Design}

An alternative to gradient-based training is to evolve the network
topology itself. The NeuroEvolution of Augmenting Topologies (NEAT)
algorithm, introduced by \citet{stanley2002evolving}, was a landmark
contribution in this direction. NEAT starts with a minimal network and
incrementally adds nodes and connections, using a principled crossover
mechanism and speciation to protect structural innovation. It
demonstrated that evolving topology alongside weights can significantly
accelerate learning on reinforcement learning benchmarks
\citep{stanley2002evolving,stanley2002efficient}.

NEAT was later extended by HyperNEAT \citep{stanley2009hypercube}, which
uses an indirect encoding called a Compositional Pattern-Producing
Network (CPPN) to generate connectivity patterns with geometric
regularity. HyperNEAT enabled the evolution of much larger networks than
direct encoding, but it still relies on a predefined substrate geometry.

A parallel line of work has explored growing neural networks with
incremental learning. The Cascade-Correlation architecture of
\citet{fahlman1990cascade} begins with a minimal network and adds hidden
units one by one, training each new unit to maximize the correlation
between its output and the residual error. The Growing Neural Gas
algorithm of \citet{fritzke1995growing} takes a different approach,
using competitive Hebbian learning to build topological representations
of input space incrementally. Both approaches demonstrate the power of
growth, but neither addresses the problem of modular structure or
recurrence.

\subsection{Modularity in Neural Networks}

Modularity—the decomposition of a system into cohesive, loosely coupled
components—is a well-established principle in software engineering
\citep{parnas1972modularity}. \citet{alho2024modularity} conducted a
systematic review of 86 studies on modular neural networks and found
that nearly two-thirds report improvements in task accuracy compared to
monolithic solutions, with 82\% reporting benefits in training and
inference time, memory, or energy consumption.

In the context of neuroevolution, \citet{khare2005coevolutionary}
proposed a co-evolutionary model for automatic problem decomposition, in
which a population of modules is evolved alongside a population of
complete systems that combine modules. They demonstrated that if a
particular task decomposition is beneficial, it can be evolved through
this mechanism. However, their model requires the designer to specify
the number of modules and the interface between them.

\citet{landassuri2012evolution} investigated the emergence of modularity
in evolved neural networks and found that pure-modular architectures
emerge at low connectivity values, and that the more different two tasks
are, the greater the modularity obtained during evolution. These results
suggest that modularity is not merely a design choice but can be an
emergent property of evolutionary search.

\subsection{Sparse Connectivity and Pruning}

Biological neural networks are highly sparse: the probability that any
two neurons are connected is on the order of $10^{-6}$ \citep{lecun2015deep}.
This sparsity is not accidental; it is a consequence of metabolic
constraints and is believed to be essential for efficient information
processing.

In artificial neural networks, sparsity has been studied primarily as a
compression technique. The Sparse Evolutionary Training (SET) method of
\citet{mocanu2018scalable} evolves an initial sparse topology into a
scale-free topology during training, allowing the training of networks
with over one million neurons on commodity hardware. More recently,
\citet{shah2026pruning} proposed an evolutionary perspective on pruning,
modeling parameter groups as populations whose influence evolves under
selection pressure.

However, these methods typically begin with a predefined network
architecture and prune or sparsify it. They do not address the question
of how sparsity and connectivity structure emerge from first principles.

\subsection{Self-Organizing and Developmental Systems}

A related line of research investigates self-organizing neural networks
inspired by biological development. \citet{plantec2024evolving} proposed
the Lifelong Neural Developmental Program (LNDP), a class of
self-organizing networks capable of synaptic and structural plasticity
in an activity- and reward-dependent manner. Their results demonstrated
that structural plasticity is advantageous in environments requiring
fast adaptation or with non-stationary rewards.

The GRN-SO-RC model of \citet{hajizada2024developmental} uses a gene
regulatory network to regulate synaptic and structural plasticity in a
reservoir computing framework, showing that self-organized reservoir
structures can adapt to specific tasks. These approaches demonstrate the
potential of developmental mechanisms, but they typically rely on
gradient-based training for the weights and a separate mechanism for
structure.

\subsection{Universality of Neural Computation}

The theoretical foundations of neural computation were established by
\citet{siegelmann1995computational}, who proved that recurrent neural
networks with rational weights are Turing complete. More recently,
\citet{perez2021turing} showed that transformers and Neural GPUs are
also Turing complete, without requiring external memory. These results
establish that neural networks can, in principle, simulate any
computable function, but they do not address the question of how such
networks can be discovered or evolved.

\section{The Present Work}

In this paper, we present \modnet{}, a self-growing modular recurrent
network that addresses the limitations of fixed architectures by
discovering its own structure through evolution. The architecture has
three defining characteristics:

\begin{enumerate}
  \item \textbf{Layer-free.} There are no predefined layers. Each node
    computes a weighted sum of its inputs followed by a threshold
    activation, and any node may connect to any other node, including
    itself.
  \item \textbf{Modular.} Nodes are organized into modules, but the
    number and size of these modules are not fixed. The network grows
    and prunes nodes within modules, and cross-module connections
    emerge as needed.
  \item \textbf{Evolved.} The topology—including the number of nodes,
    the number of modules, and all connections—is discovered through a
    compact evolutionary loop. No gradient information is used.
\end{enumerate}

We evaluate \modnet{} on a suite of six tasks spanning Boolean logic
(XOR, parity of four and six bits) and continuous regression
(sinusoids, quadratic surfaces). Our results show that \modnet{}:

\begin{itemize}
  \item Solves XOR and 4-bit parity with $100\%$ accuracy using 8 and
    12 nodes respectively;
  \item Achieves root-mean-square errors below $0.05$ on continuous
    targets with fewer than 30 nodes;
  \item Spontaneously develops cross-module connectivity (60--82\% of
    active connections), recurrence, and sparse connectivity;
  \item Exhibits two distinct growth regimes: shrinking for Boolean
    problems and growing for continuous regression.
\end{itemize}

These results suggest that layer-free computation driven by a minimal
evolutionary process can rediscover structural principles—modularity,
recurrence, sparsity—that are typically imposed by human designers.

\section{Contributions}

The main contributions of this work are:

\begin{enumerate}
  \item \textbf{Architecture.} We introduce a layer-free, modular,
    recurrent network architecture in which structure is fully
    discovered through evolution. Unlike NEAT \citep{stanley2002evolving}
    and HyperNEAT \citep{stanley2009hypercube}, our architecture does
    not assume a layered or grid-based substrate.
  \item \textbf{Emergent structure.} We demonstrate that modularity,
    cross-module integration, recurrence, and sparsity emerge
    spontaneously as the task demands, without explicit regularization
    or architectural priors.
  \item \textbf{Boolean and continuous tasks.} We show that the same
    architecture solves both discrete Boolean functions (parity) and
    continuous regression problems (sinusoids, quadratic surfaces),
    using fewer than 30 nodes in all cases.
  \item \textbf{Reproducibility.} We provide a complete, open-source
    implementation and a comprehensive experimental framework that
    generates all figures and tables reported in this paper.
\end{enumerate}

\section{Organization of the Thesis}

The remainder of this thesis is organized as follows.
Chapter~\ref{ch:related} reviews related work in neuroevolution,
modular neural networks, and self-organizing systems.
Chapter~\ref{ch:method} describes the \modnet{} architecture and the
evolutionary algorithm in detail.
Chapter~\ref{ch:experiments} presents the experimental setup, including
the tasks, the evaluation protocol, and the hyperparameters.
Chapter~\ref{ch:results} reports the results of our experiments.
Chapter~\ref{ch:discussion} discusses the implications, limitations, and
future directions.
Chapter~\ref{ch:conclusion} concludes the thesis.